\section{双积分系统}

\subsection{问题描述}

考虑 $N\ge 3$ 个机器人，编号集合为 $\mathcal{N}=\{1,2,\ldots,N\}$，每个机器人采用双积分模型
$\dot p_i=v_i,\quad \dot v_i=u_i+d_i,\quad i\in\mathcal{N},$
其中 $p_i,v_i,u_i,d_i\in\mathbb{R}^2$ 分别为位置、速度、控制输入和输入通道扰动。目标编号为 $0$，位置为 $p_0(t)\in\mathbb{R}^2$。控制目标为：在保持通信连通和避免碰撞的前提下，使所有机器人有限时间内跟踪目标。

\subsection{本算法}

整体控制方案由四部分组成：首先设计分布式有限时间观测器估计目标状态，然后构造积分型人工势场处理约束，接着定义复合误差并将其纳入 RISE 框架，最终给出完整的合围控制律。

\subsubsection{分布式有限时间观测器}

令 $\hat p_{i0}$ 和 $\hat v_{i0}$ 分别表示相对位置 $p_{i0}=p_i-p_0$ 和相对速度 $\dot p_{i0}=\dot p_i-\dot p_0$ 的估计值，定义估计误差 $e_i=\hat p_{i0}-p_{i0}$，并构造一致性变量
$\xi_i=\sum_{j\in\mathcal{N}_i(0)}a_{ij}(e_i-e_j)+b_i e_i.$
实际实现时 $\xi_i$ 可用相对测量表示为
$\xi_i=\sum_{j\in\mathcal{N}_i(0)}a_{ij}(\hat p_{i0}-\hat p_{j0}-p_{ij})+b_i(\hat p_{i0}-p_{i0}),$
其中最后一项仅对 $b_i>0$ 的机器人存在。

基于 $\xi_i$，设计如下有限时间观测器：
$$
\begin{cases}
\dot{\hat p}_{i0}=\hat v_{i0}-k_a\lfloor\xi_i\rceil^{1/2},\\
\dot{\hat v}_{i0}=u_i-k_b\operatorname{sgn}(\xi_i),
\end{cases}
\quad i\in\mathcal{N},
$$
其中 $\operatorname{sgn}(x)=x/\|x\|$（$x=0$ 时取 $\mathbf{0}$），$\lfloor x\rceil^q=\|x\|^q\operatorname{sgn}(x)$。参数选取满足
$k_b\ge \bar n/\rho,\quad k_a\ge \bigl((2\gamma_0+2\gamma_1+1)/\lambda_1\bigr)^{1/2}k_b^{1/2},$
其中 $0<\rho<1$，$\gamma_1=1+\rho$，$\gamma_0=(1+3\gamma_1)/(1-\rho)$，$\lambda_1$ 为 $\mathcal{M}(0)=\mathcal{L}(0)+\mathcal{B}$ 的最小特征值，扰动上界为 $\bar n=\max_{i\in\mathcal{N}}\bar n_i$，$\bar n_i=\varrho_2+\bar d_{0i}$。

\subsubsection{积分型人工势场}

为同时实现避碰与初始通信边保持，引入积分型人工势场（iAPF）$\phi_i=\phi_i^r+\phi_i^p$。其中避碰项为
$\phi_i^r=\sum_{j\in\mathcal{N}_i(t)}\alpha_r(\|p_{ij}\|)\frac{p_{ij}}{\|p_{ij}\|},$
函数 $\alpha_r(s)$ 在 $s\to d^+$ 时趋于 $+\infty$，在 $s\ge \mu$ 时为 $0$。连通性保持项为
$\phi_i^p=\sum_{j\in\mathcal{N}_i(0)}\alpha_p(\|p_{ij}\|)\frac{p_{ij}}{\|p_{ij}\|},$
函数 $\alpha_p(s)$ 在 $s\to \mu^-$ 时趋于 $-\infty$，在 $s\to d^+$ 时趋于 $0$。

一组可选的具体函数为
$$
\alpha_r(s)=
\begin{cases}
a_r(s)-a_r(l_d), & s\in(d,l_d),\\
0, & s\in[l_d,\infty),
\end{cases}
\quad
a_r(s)=\frac{c_r}{s-d},
$$
$$
\alpha_p(s)=
\begin{cases}
0, & s\in(d,l_\mu),\\
a_p(s)-a_p(l_\mu), & s\in[l_\mu,\mu),
\end{cases}
\quad
a_p(s)=\frac{c_p}{s-\mu},
$$
其中 $l_d,l_\mu\in(d,\mu)$，$c_r,c_p>0$。

\subsubsection{SMC 合围控制律}

引入 iAPF 辅助变量 $\Phi_{1i}$，满足 $\dot\Phi_{1i}=k_1p_{i0}-\phi_i$，$\Phi_{1i}(0)=\mathbf{0}$。在此基础上定义复合误差
$s_{i0}=\dot p_{i0}+k_2p_{i0}+\Phi_{1i}.$
该构造将相对速度误差、相对位置误差和势场补偿统一到一个变量中。由于 $p_{i0}$ 与 $\dot p_{i0}$ 不完全可测，用观测量替代得到
$\hat s_{i0}=\hat v_{i0}+k_2\hat p_{i0}+\hat\Phi_{1i},$
其中 $\hat\Phi_{1i}$ 按相同方式在线积分：$\dot{\hat\Phi}_{1i}=k_1\hat p_{i0}-\phi_i$，$\hat\Phi_{1i}(0)=\mathbf{0}$。

基于上述复合误差，设计控制输入：
$$
u_i=-k_s \mathrm{sgn}(\hat s_{i0})+\hat\Delta_i
$$
其中 $k_s>0$，$\alpha>0$，$\beta_i>0$。补偿项 $\hat\Delta_i$ 用于抵消已知的动力学分量：
$\hat\Delta_i=-k_2\hat v_{i0}-k_1\hat p_{i0}+\phi_i,$
对应的不含观测误差的真实量为 $\Delta_i=-k_2\dot p_{i0}-k_1p_{i0}+\phi_i$。
增益选取满足
$k_s \ge \sup_{t\geq 0}\|\ddot{p}_0\| + \sup_{t\geq 0,i\in\mathbb{N}}\|d_i\|$.
